
\section{Robustness Checks}
\label{section:robustness}

To verify the veracity of our findings reported above, we implement two different procedures. First, as discussed earlier, the plausibility of our identifying assumption would be in serious doubt if job market outcomes varied by ethnicity in Addis Ababa, where we expect no meaningful correlation between MTI take-up and employment, since students in the capital are equally proficient in the official working language of the city (Amharic) regardless of their ethnic backgrounds. The finding in Table \ref{tab:aareg} that being a member of the historically dominant linguistic group in Ethiopia does not confer clear advantages in educational and employment outcomes in the capital supports the validity of the identification strategy. 

These results may seem counter-intuitive, since dominant groups typically enjoy enhanced Socioeconomic status. Ethiopia's peculiar history might explain why the finding is at variance with the general principle: Ethiopia is the only African country that managed to successfully resist external colonial aggression in unison, despite facing serious internal squabbles and conflicts between its various ethnic groups, particularly between the Amhara, the Oromo, and the Tigray groups. Due to these unique circumstances, the country's ruling political and military elites have been diverse in terms of their ethnic backgrounds and have benefited from institutionalizing Amharic as the country's leading language. In other words, entry into the economic and social mainstream in the capital was available to those who subscribed to the official rules that elevated Amharic as the country's lingua franca. Hence, despite having different ethnic backgrounds and native languages, the residents of Addis Ababa are equally proficient in Amharic, and to the extent that proficiency in the language matters for success in the labor market, the finding that ethnicity has no effect on educational and employment outcomes in the capital is not unreasonable, though atypical when viewed from a general perspective.

As an additional test of robustness check, we compare educational and job market outcomes in states that have adopted the 1994 language policy similarly. Only Oromia and Tigray implemented the MTI policy fully in the survey, with Tigray adopting the Sabean script (a pre-existing script that has been in use in the country for Millennia), and Oromia introducing a new writing system based on the Roman script that has attracted strong advocates and detractors \citep{chicoine_schooling_2019}. We will investigate whether unobserved heterogeneity contributed to differences in labor market outcomes in these two states, by estimating  $\theta\ $ in a regression of \eqref{eq:03}, where all variables are as defined above except $ T_{s} $, which is now modeled as an indicator variable that is switched on for one of the states, Tigray.  

\begin{equation}\label{eq:03}
	Y_{is} = \alpha + \theta\cdot T_{s} + \beta\cdot MTIuse_{is} + X_{is}'\delta + \epsilon_{is}
\end{equation}


The estimates in Table \ref{tab:oromtig} suggest that the unobserved differences in institutions and policies, which might have not been explicitly accounted for in the empirical model, do not appear to explain the variations in employment and educational outcomes in contemporary Ethiopia.  Finding that Tigray and Oromia, which can be differentiated from each other in terms of some of their institutions and policies, do not exhibit measurable differences in their educational and employment outcomes further corroborates the key conclusions of the study. 












